\chapter{Conclusion}\label{chap:conclusion}

The systems we want to make private already have issuers, auditors,
registration authorities and central banks in them, and these entities exist
because the systems need them. In this thesis, we take their roles, and the
auditing and certification requirements that come with them, as given, and
ask how narrowly each capability can be defined. In each of our protocols,
these entities keep the functions they already have, such as providing
availability, maintaining state, certifying information or investigating
transactions, while cryptography limits what each of them can learn and do.
Users keep their privacy, and in some cases gain privacy compared with existing
designs, without taking on the work of these entities themselves.

Each chapter gives a role only the capability it needs. In \Cref{chap:tsps},
certification is split between several signers, and only a number of signers
above the threshold can sign together. In \Cref{chap:agg-creds}, the
registration authority binds each user's key to a tag and sees no attribute values, and each issuer learns only
the attributes it certifies. In \Cref{chap:txid}, auditors trace only the users
and epochs they have been granted, and a grant never exposes users who are not
targeted. In \Cref{chap:offline-payments}, withdrawals cannot be linked to
deposits, the secure element holds only the secret key and produces its part of
a joint proof, and the auditor revokes the anonymity of a double spender and no
one else.

Our constructions also avoid interaction and per-user state wherever possible.
As discussed in \Cref{chap:introduction}, issuers in different countries have
no reason to take part in an expensive interactive protocol for the sake of
users' privacy. In \Cref{chap:tsps}, all interaction between signers happens in
key generation and in an offline phase, before the message is known, and each
signer then signs locally. Issuers of aggregate credentials never coordinate
with each other and keep no state about users. Users of the tracing protocol
create transactions concurrently and in any order, without counters. Offline
payments need no network at all, since payees accept a payment from the
transcript alone and can pay the token onward.

Every chapter also gives a generic construction from standard primitives, and
then an efficient instantiation, which we implement and benchmark. Most instantiations
use Sigma protocols and pairing product equations, and the tracing protocol is
also instantiated with MiMC and a zk-SNARK. The online phase of threshold
signing stays below 2~ms up to sixteen parties. With 128 issuers and 64
disclosed attributes, an aggregate credential presentation takes 20.1~ms to
produce and 31.1~ms to verify. A smart card authorises an offline payment in
under 400~ms.

Decentralisation and privacy can feel so aligned that the involvement of any
central party in a protocol can seem like a threat to users' privacy. But a
decentralised system can reveal everything, and a central party can learn
nothing. When credential issuance is distributed across many issuers, every
issuer can learn every attribute. In \Cref{chap:agg-creds}, a single
registration authority that learns nothing about attributes keeps issuers
independent and makes aggregation cheap. Availability does not need many
servers either. The payments of \Cref{chap:offline-payments} keep working with
no server online, because the secure element prevents double spending and the
payee checks the transcript itself. We should separate the properties we want,
such as privacy, availability and accountability, from the architectures we
often reach for to achieve them.

Our results have limitations. The efficient instantiations need stronger
assumptions than the generic constructions: the GGM for $\AGHO$ in
\Cref{chap:tsps,chap:agg-creds}, $q$-type assumptions for BBS+ and the
Dodis-Yampolskiy PRF in \Cref{chap:txid}, and the DL, DDH, $q$-SDH and $q$-MSDH
assumptions in the random oracle model in \Cref{chap:offline-payments}.
Threshold $\KPW$ avoids the GGM and rests on the SXDH assumption, with
signatures of seven group elements instead of three. Our security proofs are
in the standalone model, since extraction rewinds, and composable guarantees
would need different proof techniques. The chapters share a credential layer, a
registration authority with BBS+ or SPS credentials, but we give no joint
analysis of credentials, tracing and payments run together. Offline payments
also rely on a trusted auditor to identify double spenders.

Several problems remain open. The first is weaker assumptions. A rerandomisable
SPS scheme under standard assumptions would remove the GGM from aggregate
credentials, but with larger registration signatures, and it is open
whether tracing can be built from a primitive that avoids $q$-type assumptions.
The second is stronger security models. For threshold SPS schemes, our
reduction under adaptive corruption loses a factor exponential in the number of
signers and assumes erasures, and strong unforgeability and real key generation
remain to be proven. A concrete security analysis would give the security of
aggregate credentials at the parameters we use, and composable security is
open throughout. The third is spreading trust further. The threshold techniques
of \Cref{chap:tsps} could split the auditor of \Cref{chap:txid} between several
entities, and revoking the anonymity of double spenders without a trusted
auditor is open. Revocation, of credentials and of offline payments, is
limited to expiry by epochs. Finally, constant-size presentations, issuer
hiding, retrospective grants and divisible tokens would each extend one of our
protocols.

We return to the argument of the introduction. The alternative to auditable
payments is not perfect privacy, but the status quo, where banks and payment
providers see all transactions in the clear. Auditing and certification that
are defined narrowly, to one user, one epoch or one double spender, and
enforced by cryptography, are not a hole punched in a private system. They are
how privacy can improve on what exists today, in systems that keep the entities
they already rely on.
